% ============================================================ 
% INTRODUCTION
% ============================================================
\section*{Introduction}

Throughout the twentieth century, 
numerous airfoil families were introduced, including the Clark Y and the NACA 4-digit series. 
Over time, several thousand airfoils have been developed and experimentally tested. 
The Wright brothers themselves evaluated several hundred airfoils using a homemade wind tunnel, 
marking one of the earliest systematic investigations of wing profiles.

In the early years of aviation, 
aircraft wings relied heavily on external bracing systems, 
such as struts and wires, to provide structural support. 
As stronger construction materials became available, 
engineers sought to eliminate these external supports in order to reduce aerodynamic drag. 
Achieving this objective required thicker airfoil sections capable of providing sufficient structural strength, 
leading to the evolution of airfoil geometries similar to the Clark Y profile.

Early textbooks on flight theory frequently used the Clark Y airfoil to illustrate fundamental aerodynamic principles, and this practice has persisted over time. 
However, the airfoil families developed by the National Advisory Committee for Aeronautics (NACA) were documented more systematically. 
In particular, NACA Report No.~460 \cite{NACA_TR_460}, 
published in November 1933 and entitled \textit{The Characteristics of the Seventy-Eight Related Airfoil Sections from Tests in the Variable Density Wind Tunnel}, 
provided detailed experimental data for a series of related airfoils.

As a governmental research agency, 
NACA conducted extensive wind tunnel testing to determine the aerodynamic efficiency of various airfoil configurations under different flight conditions. 
NACA later evolved into the National Aeronautics and Space Administration (NASA), which continued research and development in airfoil design. 
The resulting NACA airfoil profiles were designated using a standardized numbering system 
consisting of four, five, or six digits, enabling a direct relationship between numerical parameters and geometric characteristics.